\documentclass[11pt]{article}

\usepackage[a4paper,margin=1in]{geometry}
\usepackage[T1]{fontenc}
\usepackage{lmodern}
\usepackage{amsmath,amssymb}
\usepackage{microtype}
\usepackage{caption}
\usepackage{tikz}
\usetikzlibrary{
    positioning,
    calc,
    fit,
    backgrounds,
    arrows.meta,
    decorations.pathreplacing,
    patterns,
    matrix,
    shapes.geometric,
    shapes.misc
}

\title{Monochrome TikZ Figure Sheet for Machine Learning and Deep Learning}
\author{}
\date{}

% =========================================================
% Monochrome TikZ Style Library
% =========================================================
\tikzset{
    >=Stealth,
    every node/.style={font=\small},

    % ---------- Semantic node styles ----------
    bwmodule/.style={
        draw=black,
        thick,
        rounded corners=2pt,
        fill=white,
        minimum width=28mm,
        minimum height=9mm,
        align=center
    },
    bwprocess/.style={
        draw=black,
        thick,
        rectangle,
        fill=white,
        minimum width=22mm,
        minimum height=8mm,
        align=center
    },
    bwdata/.style={
        draw=black,
        thin,
        rectangle,
        fill=gray!12,
        minimum width=22mm,
        minimum height=8mm,
        align=center
    },
    bwlatent/.style={
        draw=black,
        thick,
        dashed,
        rectangle,
        fill=white,
        minimum width=22mm,
        minimum height=8mm,
        align=center
    },
    bwnoise/.style={
        draw=black,
        thick,
        rectangle,
        pattern=dots,
        pattern color=black,
        minimum width=22mm,
        minimum height=8mm,
        align=center
    },
    bwop/.style={
        draw=black,
        thick,
        circle,
        fill=white,
        minimum size=6mm,
        inner sep=0pt
    },
    bwnote/.style={
        draw=black,
        thin,
        dashed,
        rounded corners=2pt,
        fill=white,
        minimum width=24mm,
        minimum height=8mm,
        align=center
    },

    % ---------- Arrows ----------
    bwmain/.style={->, thick, black},
    bwaux/.style={->, thin, black},
    bwskip/.style={->, dashed, thick, black},
    bwconcept/.style={->, dotted, thick, black},
    bwinverse/.style={->, very thick, black},

    % ---------- Group / repetition ----------
    bwgroup/.style={
        draw=black,
        thick,
        dashed,
        rounded corners=3pt,
        inner sep=6pt
    },
    bwbrace/.style={
        decorate,
        decoration={brace, amplitude=4pt},
        thick
    },

    % ---------- Annotation ----------
    bwannot/.style={
        font=\scriptsize,
        align=center
    }
}

% ---------- Convenience macros ----------
\newcommand{\bwmodule}[4][]{\node[bwmodule,#1] (#2) at #3 {#4};}
\newcommand{\bwprocess}[4][]{\node[bwprocess,#1] (#2) at #3 {#4};}
\newcommand{\bwdata}[4][]{\node[bwdata,#1] (#2) at #3 {#4};}
\newcommand{\bwlatent}[4][]{\node[bwlatent,#1] (#2) at #3 {#4};}
\newcommand{\bwnoise}[4][]{\node[bwnoise,#1] (#2) at #3 {#4};}
\newcommand{\bwnote}[4][]{\node[bwnote,#1] (#2) at #3 {#4};}

\newcommand{\bwconnect}[2]{\draw[bwmain] (#1) -- (#2);}
\newcommand{\bwauxconnect}[2]{\draw[bwaux] (#1) -- (#2);}
\newcommand{\bwskipconnect}[2]{\draw[bwskip] (#1) -- (#2);}
\newcommand{\bwconceptconnect}[2]{\draw[bwconcept] (#1) -- (#2);}

\begin{document}

\maketitle

\section*{Figure 1: Transformer Architecture}

\begin{figure}[h!]
\centering
\begin{tikzpicture}[x=1cm,y=1cm]

% ---------- Encoder input ----------
\bwdata{xin}{(0,0)}{Input Tokens\\$x_1,\dots,x_n$}
\bwprocess{embin}{(0,1.5)}{Token\\Embedding}
\bwprocess{posin}{(0,3.0)}{+ Positional\\Encoding}

\bwconnect{xin}{embin}
\bwconnect{embin}{posin}

% ---------- Encoder stack ----------
\bwmodule{encattn}{(0,5.0)}{Multi-Head\\Self-Attention}
\bwprocess{encnorm1}{(0,6.6)}{Add \& Norm}
\bwmodule{encffn}{(0,8.2)}{Feed Forward}
\bwprocess{encnorm2}{(0,9.8)}{Add \& Norm}

\bwconnect{posin}{encattn}
\bwconnect{encattn}{encnorm1}
\bwconnect{encnorm1}{encffn}
\bwconnect{encffn}{encnorm2}

\draw[bwskip] ($(posin.east)+(0.2,0)$) -- ++(0.9,0) |- (encnorm1.east);
\draw[bwskip] ($(encnorm1.east)+(0.2,0)$) -- ++(0.9,0) |- (encnorm2.east);

\begin{scope}[on background layer]
\node[bwgroup, fit=(encattn)(encnorm1)(encffn)(encnorm2), inner sep=7pt] (encbox) {};
\end{scope}

\draw[bwbrace] (1.8,4.3) -- (1.8,10.5);
\node[bwannot, right=4pt] at (1.8,7.4) {Encoder block\\repeated $N$ times};

% ---------- Decoder input ----------
\bwdata{yshift}{(7,0)}{Shifted Output\\Tokens}
\bwprocess{embout}{(7,1.5)}{Token\\Embedding}
\bwprocess{posout}{(7,3.0)}{+ Positional\\Encoding}

\bwconnect{yshift}{embout}
\bwconnect{embout}{posout}

% ---------- Decoder stack ----------
\bwmodule{decself}{(7,5.0)}{Masked Multi-Head\\Self-Attention}
\bwprocess{decnorm1}{(7,6.6)}{Add \& Norm}
\bwmodule{crossattn}{(7,8.2)}{Encoder--Decoder\\Attention}
\bwprocess{decnorm2}{(7,9.8)}{Add \& Norm}
\bwmodule{decffn}{(7,11.4)}{Feed Forward}
\bwprocess{decnorm3}{(7,13.0)}{Add \& Norm}

\bwconnect{posout}{decself}
\bwconnect{decself}{decnorm1}
\bwconnect{decnorm1}{crossattn}
\bwconnect{crossattn}{decnorm2}
\bwconnect{decnorm2}{decffn}
\bwconnect{decffn}{decnorm3}

\draw[bwskip] ($(posout.east)+(0.2,0)$) -- ++(0.9,0) |- (decnorm1.east);
\draw[bwskip] ($(decnorm1.east)+(0.2,0)$) -- ++(0.9,0) |- (decnorm2.east);
\draw[bwskip] ($(decnorm2.east)+(0.2,0)$) -- ++(0.9,0) |- (decnorm3.east);

\begin{scope}[on background layer]
\node[bwgroup, fit=(decself)(decnorm1)(crossattn)(decnorm2)(decffn)(decnorm3), inner sep=7pt] (decbox) {};
\end{scope}

\draw[bwbrace] (8.9,4.3) -- (8.9,13.7);
\node[bwannot, right=4pt] at (8.9,9.0) {Decoder block\\repeated $N$ times};

% ---------- Cross attention memory ----------
\draw[bwmain] (encnorm2.east) -- ++(1.2,0) |- (crossattn.west);
\node[bwannot] at (4.0,9.1) {encoder memory};

% ---------- Output head ----------
\bwprocess{lin}{(7,14.8)}{Linear}
\bwprocess{softmax}{(7,16.3)}{Softmax}
\bwdata{probs}{(7,17.8)}{Next-token\\distribution}

\bwconnect{decnorm3}{lin}
\bwconnect{lin}{softmax}
\bwconnect{softmax}{probs}

% ---------- Side labels ----------
\node[bwannot] at (-2.0,5.0) {self-attention over\\input sequence};
\node[bwannot] at (10.8,5.0) {causal masking};
\node[bwannot] at (10.9,8.2) {attention to encoder\\representations};

\end{tikzpicture}
\caption{Monochrome Transformer architecture. Rounded rectangles denote learned modules, sharp rectangles denote deterministic operations, gray boxes denote token/data objects, and dashed arrows denote residual connections.}
\end{figure}

\clearpage

\section*{Figure 2: Diffusion Model Pipeline}

\begin{figure}[h!]
\centering
\begin{tikzpicture}[x=1cm,y=1cm]

% ---------- Forward noising ----------
\bwdata{x0}{(0,0)}{$x_0$\\clean data}
\bwdata{x1}{(3,0)}{$x_1$}
\bwdata{xt}{(6,0)}{$x_t$}
\bwnoise{xT}{(9,0)}{$x_T$\\pure noise}

\draw[bwmain] (x0) -- node[above, bwannot] {add noise} (x1);
\draw[bwmain] (x1) -- node[above, bwannot] {$\cdots$} (xt);
\draw[bwmain] (xt) -- node[above, bwannot] {repeat} (xT);

\node[bwannot] at (4.5,1.2) {Forward diffusion\\$q(x_t\mid x_{t-1})$};

% ---------- Reverse denoising ----------
\bwnoise{zT}{(9,-4)}{$x_T$}
\bwdata{zt}{(6,-4)}{$x_t$}
\bwdata{ztm}{(3,-4)}{$x_{t-1}$}
\bwdata{xhat}{(0,-4)}{$x_0$\\generated sample}

\draw[bwinverse] (zT) -- node[below, bwannot] {denoise} (zt);
\draw[bwinverse] (zt) -- node[below, bwannot] {$\cdots$} (ztm);
\draw[bwinverse] (ztm) -- node[below, bwannot] {repeat} (xhat);

\node[bwannot] at (4.5,-5.2) {Reverse process\\$p_\theta(x_{t-1}\mid x_t)$};

% ---------- Denoiser network ----------
\bwmodule{denoiser}{(6,-1.9)}{Noise predictor\\$\epsilon_\theta(x_t,t)$}
\draw[bwaux] (xt) -- (denoiser);
\draw[bwconcept] (denoiser) -- (zt);

% ---------- Time embedding ----------
\bwprocess{time}{(10.7,-1.9)}{time\\embedding $t$}
\draw[bwaux] (time.west) -- (denoiser.east);

% ---------- Closed-form forward note ----------
\bwnote{cf}{(2.7,2.8)}{Closed-form marginal\\
$x_t=\sqrt{\bar\alpha_t}\,x_0+\sqrt{1-\bar\alpha_t}\,\epsilon$}
\draw[bwconcept] (cf.south) -- (xt.north);

% ---------- Training objective ----------
\bwnote{loss}{(2.8,-6.8)}{Training objective\\
$\mathcal L=\mathbb E\|\epsilon-\epsilon_\theta(x_t,t)\|^2$}
\draw[bwconcept] (denoiser.south west) to[out=-130,in=20] (loss.north east);

% ---------- Labels ----------
\node[bwannot] at (10.5,0.9) {increasing noise level};
\node[bwannot] at (10.7,-4.9) {iterative refinement};

\end{tikzpicture}
\caption{Monochrome diffusion pipeline. Gray boxes denote data states, the dotted-pattern box denotes pure noise, and the lower chain represents iterative denoising.}
\end{figure}

\clearpage

\section*{Figure 3: Normalizing Flow}

\begin{figure}[h!]
\centering
\begin{tikzpicture}[x=1cm,y=1cm]

% ---------- Endpoints ----------
\bwlatent{z}{(0,0)}{Latent variable\\$z\sim p_Z(z)$}
\bwdata{x}{(12,0)}{Data sample\\$x\sim p_X(x)$}

% ---------- Flow blocks ----------
\bwmodule{f1}{(2.5,0)}{$f_1$}
\bwmodule{f2}{(5.0,0)}{$f_2$}
\bwmodule{f3}{(7.5,0)}{$\cdots$}
\bwmodule{f4}{(10.0,0)}{$f_K$}

\draw[bwmain] (z) -- (f1);
\draw[bwmain] (f1) -- (f2);
\draw[bwmain] (f2) -- (f3);
\draw[bwmain] (f3) -- (f4);
\draw[bwmain] (f4) -- (x);

\node[bwannot] at (6,1.1) {Forward map: $x=f(z)$};

% ---------- Intermediate states ----------
\node[bwannot] at (2.5,-1.0) {$h_1$};
\node[bwannot] at (5.0,-1.0) {$h_2$};
\node[bwannot] at (7.5,-1.0) {$\cdots$};
\node[bwannot] at (10.0,-1.0) {$h_K$};

% ---------- Inverse direction ----------
\draw[bwinverse] (x.south) to[out=-90,in=-90] node[below, bwannot] {inverse map $z=f^{-1}(x)$} (z.south);

% ---------- Notes ----------
\bwnote{simple}{(0,2.8)}{Simple base\\distribution}
\bwnote{complex}{(12,2.8)}{Complex target\\distribution}
\bwnote{invnote}{(6,2.8)}{Each block is invertible\\with tractable Jacobian determinant}

\draw[bwconcept] (simple.south) -- (z.north);
\draw[bwconcept] (complex.south) -- (x.north);
\draw[bwconcept] (invnote.south) -- (f2.north);
\draw[bwconcept] (invnote.south) -- (f4.north);

% ---------- Change-of-variables ----------
\bwnote[minimum width=60mm, minimum height=14mm]{cov}{(6,-3.3)}{Change of variables:
\[
p_X(x)=p_Z(z)\left|\det\frac{\partial f^{-1}(x)}{\partial x}\right|
\]
Layerwise:
\[
\log p_X(x)=\log p_Z(z)+\sum_{k=1}^K
\log\left|\det\frac{\partial h_{k-1}}{\partial h_k}\right|
\]
};

\end{tikzpicture}
\caption{Monochrome normalizing flow diagram. Dashed boxes denote conceptual latent objects, gray boxes denote observed data, and rounded rectangles denote invertible learned transformations.}
\end{figure}

\clearpage

\section*{Figure 4: Affine Coupling Layer}

\begin{figure}[h!]
\centering
\begin{tikzpicture}[x=1cm,y=1cm]

\bwdata{xA}{(0,1.2)}{$x_A$}
\bwdata{xB}{(0,-1.2)}{$x_B$}

\bwmodule{snet}{(3,1.2)}{$s(x_A)$}
\bwmodule{tnet}{(3,-1.2)}{$t(x_A)$}

\bwprocess{idA}{(6,1.2)}{identity}
\bwprocess{affB}{(6,-1.2)}{$x_B' = x_B\odot e^{s(x_A)} + t(x_A)$}

\bwdata{yA}{(9,1.2)}{$y_A$}
\bwdata{yB}{(9,-1.2)}{$y_B$}

\draw[bwmain] (xA) -- (snet);
\draw[bwmain] (xA) -- (tnet);
\draw[bwmain] (xA) -- (idA);
\draw[bwmain] (idA) -- (yA);

\draw[bwmain] (xB) -- (affB);
\draw[bwaux] (snet) -- (affB);
\draw[bwaux] (tnet) -- (affB);
\draw[bwmain] (affB) -- (yB);

\node[bwannot] at (4.6,-2.5) {One part remains unchanged; the other is transformed conditionally.};

\end{tikzpicture}
\caption{Monochrome affine coupling layer used in normalizing flows such as RealNVP.}
\end{figure}

\end{document}